\scanpage{89}
and that $\inserted{G\mapsto}S\Pi M_G$ is a functor in \deleted{$\unsure{M}_G$} $G$ for arbitrary morphisms, and that (124) is functorial for isomorphisms. For an arbitrary morphism $G\xrightarrow{f}G'$, corresponding to a subset $C\subset A_G$ (such that $G\unsure{\rightharpoondown|{-}}G'$ is isomorphic, as a quotient of $G$, to $G\unsure{\rightharpoondown|{-}}G/C$), we assume that we have a commutative diagram
\[
\tag{125}
\diagram{\begin{tikzcd}[ampersand replacement=\&, column sep=large,row sep=large]
S\Pi M_G \arrow[r,"S\Pi(f)"] \arrow[d,"\varphi_G"] \& S\Pi M_{G'} \arrow[d,"\varphi_{G'}"] \\
S_{\deleted{\unsure{A}}C}M_G \arrow[r,"\gamma_f"] \& \deleted{\unsure{S}}\Pi M_{G'}
\end{tikzcd}}
\]
i.e. that the composite
\[
\diagram{\begin{tikzcd}[ampersand replacement=\&, column sep=large]
S\Pi M_G \arrow[r,"S\Pi(f)"] \& S\Pi M_{G'} \arrow[r,phantom,"{\overset{\varphi_{G'}}{\unsure{\rightharpoondown|{-}}}}" description] \& \Pi M_{G'}
\end{tikzcd}}
\]
factors through the extension $S_{\deleted{\unsure{A}}C}M_G$ of $\Pi M_G$ by $\mathbb Z^C$, deduced from the extension $S\Pi M_G$ of $\Pi M_G$ by $\mathbb Z^{A(G)}$ by means of the canonical homomorphism
\[
\diagram{\begin{tikzcd}[ampersand replacement=\&]
\mathbb Z^{A(G)} \arrow[r] \& \mathbb Z^C.
\end{tikzcd}}
\]
One then automatically has the transitivity relation for a composite.
\margin{88}
